% GENERATED FILE. Edit canonical lessons, then run npm run generate:books.
% canonical-source-hash: fnv1a64:42a2446024dfef1f
% canonical-lessons: AR-C531-murahiq, AR-C531-kabira, AR-C531-nadija, AR-C531-shaykhukha, AR-C531-murabbiya

\chapter{Growing up and growing old}
\label{ch:ar-growing-up-and-growing-old}
\hlchaptermodality{531}

\begin{chapteropening}
\textbf{By the end of this chapter:} \emph{I can name a teenager, a nanny and old age, and say that someone got older or matured.}

Complete the last lesson of chapter 531: I can name a teenager, a nanny and old age, and say that someone got older or matured.
\end{chapteropening}

\section[\texorpdfstring{\ar{مراهق} (murāhiq)}{مراهق (murāhiq)}]{\ar{مراهق} (murāhiq) — a teenager}

\label{lesson:AR-C531-murahiq}



\begin{quote}
Before the new one: say the Arabic for to wean, then the Arabic for to bring up.
\end{quote}

\subsection*{You'll want to know: \ar{مراهق}}
\textbf{\ar{مراهق}} — \emph{murāhiq} — \textquotedblleft{}a teenager\textquotedblright{}.

\begin{quote}
\ar{ابن} \ar{جارنا} \ar{مراهق} \ar{ويحب} \ar{الموسيقى} \ar{العالية}.
\end{quote}

\begin{quote}
\emph{ibn jārinā murāhiq wa-yuḥibbu al-mūsīqā al-ʿāliya.}
\end{quote}

\begin{quote}
Our neighbour's son is a teenager and loves loud music.
\end{quote}

\begin{cousinweb}
\textbf{\ar{مراهق}} (\emph{murāhiq}) is built on the native Arabic root \emph{r-h-q}, \textquotedblleft{}to come near\textquotedblright{} (Lane). Murāhiq is the Form III participle: one coming near to manhood.
\end{cousinweb}

\begin{grammarlens}[title={What you've built}]
One more word for describing what you lived and felt.
\end{grammarlens}

\subsection*{Your turn}
\begin{itemize}
\raggedright
  \item \emph{Say it:} \emph{murāhiq}
  \item \emph{Say it:} \emph{murāhiq}, once more
  \item \emph{Say it:} \emph{ibn jārinā murāhiq wa-yuḥibbu al-mūsīqā al-ʿāliya.} — Our neighbour's son is a teenager and loves loud music.
  \item \emph{Recall:} say the Arabic for to wean, then the Arabic for to bring up, then say \emph{murāhiq} again
\end{itemize}

\begin{tcolorbox}[breakable,colback=teal!4,colframe=teal!35!black,title={Before you move on}]
What does \ar{مراهق} mean? (\textquotedblleft{}A teenager\textquotedblright{}.) Say it once more.
\end{tcolorbox}
\section[\texorpdfstring{\ar{كبر} (kabira)}{كبر (kabira)}]{\ar{كبر} (kabira) — to get older; to grow up}

\label{lesson:AR-C531-kabira}



\begin{quote}
Before the new one: say the Arabic for to bring up, then the Arabic for a teenager.
\end{quote}

\subsection*{You'll want to know: \ar{كبر}}
\textbf{\ar{كبر}} — \emph{kabira} — \textquotedblleft{}to get older; to grow up\textquotedblright{}.

\begin{quote}
\ar{كبر} \ar{الأطفال} \ar{وسافروا} \ar{إلى} \ar{المدينة}.
\end{quote}

\begin{quote}
\emph{kabira al-aṭfāl wa-sāfarū ilā al-madīna.}
\end{quote}

\begin{quote}
The children grew up and went off to the city.
\end{quote}

Present: \ar{يكبر}, \emph{yakbar}.

Kabira, yakbaru, is to grow older; kabura, yakburu, with the same letters, is to be big or great.

\begin{cousinweb}
\textbf{\ar{كبر}} (\emph{kabira}) is built on the native Arabic root \emph{k-b-r}, \textquotedblleft{}to be big, great\textquotedblright{} (Lane). Kabīr, big, which you know, is from the same root.
\end{cousinweb}

\begin{grammarlens}[title={What you've built}]
One more word for describing what you lived and felt.
\end{grammarlens}

\subsection*{Your turn}
\begin{itemize}
\raggedright
  \item \emph{Say it:} \emph{kabira}
  \item \emph{Say it:} \emph{kabira}, once more
  \item \emph{Say it:} \emph{kabira al-aṭfāl wa-sāfarū ilā al-madīna.} — The children grew up and went off to the city.
  \item \emph{Recall:} say the Arabic for to bring up, then the Arabic for a teenager, then say \emph{kabira} again
\end{itemize}

\begin{tcolorbox}[breakable,colback=teal!4,colframe=teal!35!black,title={Before you move on}]
What does \ar{كبر} mean? (\textquotedblleft{}To get older; to grow up\textquotedblright{}.) Say it once more.
\end{tcolorbox}
\section[\texorpdfstring{\ar{نضج} (naḍija)}{نضج (naḍija)}]{\ar{نضج} (naḍija) — to ripen; to mature}

\label{lesson:AR-C531-nadija}



\begin{quote}
Before the new one: say the Arabic for a teenager, then the Arabic for to get older.
\end{quote}

\subsection*{You'll want to know: \ar{نضج}}
\textbf{\ar{نضج}} — \emph{naḍija} — \textquotedblleft{}to ripen; to mature\textquotedblright{}.

\begin{quote}
\ar{نضج} \ar{التين} \ar{في} \ar{آخر} \ar{الصيف}.
\end{quote}

\begin{quote}
\emph{naḍija at-tīn fī ākhir aṣ-ṣayf.}
\end{quote}

\begin{quote}
The figs ripened at the end of the summer.
\end{quote}

Present: \ar{ينضج}, \emph{yanḍaj}.

Of a person it means to grow mature.

\begin{cousinweb}
\textbf{\ar{نضج}} (\emph{naḍija}) is built on the native Arabic root \emph{n-ḍ-j}, \textquotedblleft{}to be ripe, to be well cooked\textquotedblright{} (Lane). Nāḍij, ripe, which you know, is its participle.
\end{cousinweb}

\begin{grammarlens}[title={What you've built}]
One more word for describing what you lived and felt.
\end{grammarlens}

\subsection*{Your turn}
\begin{itemize}
\raggedright
  \item \emph{Say it:} \emph{naḍija}
  \item \emph{Say it:} \emph{naḍija}, once more
  \item \emph{Say it:} \emph{naḍija at-tīn fī ākhir aṣ-ṣayf.} — The figs ripened at the end of the summer.
  \item \emph{Recall:} say the Arabic for a teenager, then the Arabic for to get older, then say \emph{naḍija} again
\end{itemize}

\begin{tcolorbox}[breakable,colback=teal!4,colframe=teal!35!black,title={Before you move on}]
What does \ar{نضج} mean? (\textquotedblleft{}To ripen; to mature\textquotedblright{}.) Say it once more.
\end{tcolorbox}
\section[\texorpdfstring{\ar{شيخوخة} (shaykhūkha)}{شيخوخة (shaykhūkha)}]{\ar{شيخوخة} (shaykhūkha) — old age}

\label{lesson:AR-C531-shaykhukha}



\begin{quote}
Before the new one: say the Arabic for to get older, then the Arabic for to ripen.
\end{quote}

\subsection*{You'll want to know: \ar{شيخوخة}}
\textbf{\ar{شيخوخة}} — \emph{shaykhūkha} — \textquotedblleft{}old age\textquotedblright{}.

\begin{quote}
\ar{في} \ar{الشيخوخة} \ar{يحتاج} \ar{الإنسان} \ar{إلى} \ar{الراحة}.
\end{quote}

\begin{quote}
\emph{fī ash-shaykhūkha yaḥtāju al-insān ilā ar-rāḥa.}
\end{quote}

\begin{quote}
In old age a person needs rest.
\end{quote}

\begin{cousinweb}
\textbf{\ar{شيخوخة}} (\emph{shaykhūkha}) is built on the native Arabic root \emph{sh-y-kh}, \textquotedblleft{}to grow old\textquotedblright{} (Lane). Shaykh, an old man and an elder, is from the same root.
\end{cousinweb}

\begin{grammarlens}[title={What you've built}]
One more word for describing what you lived and felt.
\end{grammarlens}

\subsection*{Your turn}
\begin{itemize}
\raggedright
  \item \emph{Say it:} \emph{shaykhūkha}
  \item \emph{Say it:} \emph{shaykhūkha}, once more
  \item \emph{Say it:} \emph{fī ash-shaykhūkha yaḥtāju al-insān ilā ar-rāḥa.} — In old age a person needs rest.
  \item \emph{Recall:} say the Arabic for to get older, then the Arabic for to ripen, then say \emph{shaykhūkha} again
\end{itemize}

\begin{tcolorbox}[breakable,colback=teal!4,colframe=teal!35!black,title={Before you move on}]
What does \ar{شيخوخة} mean? (\textquotedblleft{}Old age\textquotedblright{}.) Say it once more.
\end{tcolorbox}
\section[\texorpdfstring{\ar{مربية} (murabbiya)}{مربية (murabbiya)}]{\ar{مربية} (murabbiya) — a nanny}

\label{lesson:AR-C531-murabbiya}



\begin{quote}
Before the new one: say the Arabic for to ripen, then the Arabic for old age.
\end{quote}

\subsection*{You'll want to know: \ar{مربية}}
\textbf{\ar{مربية}} — \emph{murabbiya} — \textquotedblleft{}a nanny\textquotedblright{}.

\begin{quote}
\ar{تعتني} \ar{المربية} \ar{بالأطفال} \ar{في} \ar{النهار}.
\end{quote}

\begin{quote}
\emph{taʿtanī al-murabbiya bi-l-aṭfāl fī an-nahār.}
\end{quote}

\begin{quote}
The nanny looks after the children during the day.
\end{quote}

\begin{cousinweb}
\textbf{\ar{مربية}} (\emph{murabbiya}) is built on the native Arabic root \emph{r-b-w}, \textquotedblleft{}to grow, to increase\textquotedblright{} (Lane). Murabbiya is the feminine participle of rabbā, to bring up.
\end{cousinweb}

\begin{grammarlens}[title={What you've built}]
One more word for describing what you lived and felt.
\end{grammarlens}

\subsection*{Your turn}
\begin{itemize}
\raggedright
  \item \emph{Say it:} \emph{murabbiya}
  \item \emph{Say it:} \emph{murabbiya}, once more
  \item \emph{Say it:} \emph{taʿtanī al-murabbiya bi-l-aṭfāl fī an-nahār.} — The nanny looks after the children during the day.
  \item \emph{Recall:} say the Arabic for to ripen, then the Arabic for old age, then say \emph{murabbiya} again
\end{itemize}

\begin{tcolorbox}[breakable,colback=teal!4,colframe=teal!35!black,title={Before you move on}]
What does \ar{مربية} mean? (\textquotedblleft{}A nanny\textquotedblright{}.) Say it once more.
\end{tcolorbox}
